% appA-prelims.tex -- Additional preliminaries.
\section{Additional preliminaries}\label{app:prelims}

\paragraph{Negative association: the closure properties we use.}
We use two standard closure properties of negatively associated (NA) families, both from the original treatment of
\cite{JoagDevProschan1983}, together with the one tail inequality they imply. Recall that a random vector is NA if, for every
pair of disjoint index sets and every pair of coordinatewise-nondecreasing functions of the corresponding subvectors, the two
resulting functions are nonpositively correlated. The two closure properties we need are the following.

The first (property P7, closure under independent union) says that NA is preserved by concatenating independent NA vectors:
if $X$ and $Y$ are each NA and are independent of each other, then the stacked vector $(X,Y)$ is NA. The reason is that any
two monotone functions on disjoint index sets either read coordinates of a single block, where the NA of that block applies,
or split across the two blocks, where independence makes the covariance vanish; in every case the defining inequality holds.

The second (property P6, monotone functions of disjoint blocks) says that NA is preserved under taking nondecreasing
functions of disjoint groups of coordinates: if $(Z_1,\dots,Z_m)$ is NA and $g_1,\dots,g_r$ are coordinatewise-nondecreasing
functions, each depending only on its own \emph{disjoint} block of indices, then $(g_1(Z),\dots,g_r(Z))$ is again NA. This is
immediate from the definition, since any monotone function of a subcollection of the $g_i$ is itself a monotone function of a
union of the underlying disjoint blocks.

These two facts are exactly what is needed to certify that the row loads of a random column-sampled matrix are NA. Fix a
subset $S$ of $s$ columns. The columns are sampled independently, and within a single column the entries form an NA vector
(a vector of indicators of a uniform $d$-subset is NA by the permutation-distribution case of \cite{JoagDevProschan1983}).
Applying (P7) across the $s$ independent columns of $S$ therefore yields an NA vector on the product index set
$[n]\times S$. Now apply (P6) with the functions $g_i(\cdot)=\sum_{t\in S}A_{it}$, one per row $i$. The function $g_i$ reads
only the block $\{(i,t):t\in S\}$, these blocks are disjoint across the rows $i$, and each $g_i$ is nondecreasing in its
entries, so the hypotheses of (P6) hold and the row-load vector $\{r_i(S)\}_{i\in[n]}$ is NA.

The single inequality we actually invoke downstream is the upper-tail (MGF) direction of NA: for nondecreasing functions
$f_1,\dots,f_n$ and any $t>0$,
\[
  \E\Big[\textstyle\prod_i f_i(Y_i)\Big]\ \le\ \prod_i\E[f_i(Y_i)],
\]
which is the statement that the moment generating function of a sum of NA coordinates factorizes as an upper bound. It is the
direction of NA that makes Chernoff-type bounds go through as if the coordinates were independent. We must check that we apply
it with the inequality pointing the right way, that is, with genuinely nondecreasing $f_i$. In both estimates of
\Cref{app:structural-full} the summand $Y_i$ is a nondecreasing function of the load $r_i$: for the dense estimate it is
$\exp(-c/r)\1[r>0]$, which increases as $r$ increases (a larger load makes the exponential closer to $1$ and keeps the
indicator on); for the sparse estimate it is the threshold indicator $\1[r>U]$, which is also nondecreasing in $r$. Hence in
each case $e^{tY_i}$ is a nondecreasing function of $Y_i$ and so of $r_i$, the factorization above applies in the correct
(upper-bound) direction, and the moment generating function of $\sum_i Y_i$ is dominated by the product of the per-coordinate
moment generating functions.

\paragraph{The exact-equality form of Lovett--Meka.}
In \Cref{lem:LM} a threshold $c_j=0$ on a vector $v_j$ contributes $e^{-0}=1$ to the budget $\sum_j e^{-c_j^2/16}\le m/16$ and
enforces $\langle v_j,x-x_0\rangle=0$ exactly. We use this only for $v_j=\mathbf 1$ (the all-ones vector), so the budget
spent is $1$; together with the row constraints ($\le m/64$) it stays under $m/16$ as soon as $m\ge 22$. The partial-coloring
walk is run as a Brownian motion constrained to the subspace $\{x:\langle\mathbf 1,x\rangle=\langle\mathbf 1,x_0\rangle\}$,
the standard way to impose an exact linear equality; the cardinality of the output is therefore exact up to the $\Oh(1)$
coordinates left fractional when the alive set drops below $22$, which are rounded by hand at additive cost $\Oh(1)$.

\paragraph{Poisson tail, full statement.}
We record the upper tail of a Poisson variable in the Bennett form used in \eqref{eq:poisson-tail}, together with its
derivation, its two regimes, and the transfer to the binomial. Let $X\sim\Pois(\mu)$ and let $t\ge0$. For any $\theta>0$ the
Chernoff bound gives
\[
  \Pr(X\ge\mu+t)\ \le\ e^{-\theta(\mu+t)}\,\E[e^{\theta X}]\ =\ \exp\!\big(\mu(e^{\theta}-1)-\theta(\mu+t)\big),
\]
using the Poisson moment generating function $\E[e^{\theta X}]=\exp(\mu(e^{\theta}-1))$. Minimizing the exponent over
$\theta$ sets $e^{\theta}=1+t/\mu$, that is $\theta=\log(1+t/\mu)$, and substituting this optimal $\theta$ collapses the
exponent to $-\mu\,h(t/\mu)$ with
\[
  h(u)\ =\ (1+u)\log(1+u)-u\ \ge\ 0 .
\]
Thus $\Pr(X\ge\mu+t)\le\exp(-\mu\,h(t/\mu))$, which is the upper bound quoted in \eqref{eq:poisson-tail}. The function $h$ is
nonnegative because $h(0)=0$ and $h'(u)=\log(1+u)\ge0$, so it is increasing on $[0,\infty)$. A matching lower bound of the
same exponential order holds for $\mu\ge1$ and $t\le\mu^{2}$ by the local central limit theorem, which gives the Poisson
point mass at $\mu+t$ up to polynomial factors and lets one sum the tail with no loss in the exponent; this is what makes the
estimate tight and not merely an upper bound.

The two regimes of \eqref{eq:poisson-tail} are read off from the behavior of $h$. For small deviations, $u\le1$, a Taylor
expansion gives $h(u)=\tfrac12 u^2-\tfrac16u^3+\cdots\asymp u^2$, so the bound is the Gaussian one
$\exp(-\Th(t^2/\mu))$. For large deviations, $u\ge1$, the leading term is $h(u)\asymp u\log(1+u)$, so the bound becomes the
Poissonian $\exp(-\Th(t\log(1+t/\mu)))$, which decays faster than any Gaussian. The crossover at $u\asymp1$, that is at
$t\asymp\mu$, is exactly the point where the deviation stops being typical-scale and the heavier Poisson tail takes over.

The same estimate transfers verbatim to the binomial. For $X\sim\Bin(m,p)$ with mean $\mu=mp$, Bennett's inequality gives the
identical $\exp(-\mu\,h(t/\mu))$ upper tail, and the matching lower bound of the same order holds once the variance is at
least constant, $mp(1-p)\ge1$, so that the local central limit theorem applies. This variance condition holds throughout our
application, because the relevant binomial means are $\Th(\lam)=\Om(1)$ in the regime $k\le d$, and the success probabilities
are bounded away from $1$, so $mp(1-p)=\Th(\lam)\ge1$.
